\subsection{Fault injection}
In computer science, many types of faults can cause a system to behave unexpectedly.
Faults may occur during the development or execution of a program and can originate from the software, the hardware, or both.
Faults come in a wide variety: they can be syntax errors that cause the script to fail to compile, generated warnings, or crashes with an error during compilation.
Faults can be caused by human error, misconfigurations that were not anticipated, or malicious individuals exploiting security vulnerabilities.
One important type of fault is \textbf{fault injection}.
Fault injection has been around for a long time, and we found traces of papers dating from the 1970s, but they were never published; the oldest one published is from 1995 \cite{1_Segall1995}.
This fault can be caused by various techniques, including hardware, such as clock glitches \cite{15obermaier2017fuzzy,16korak2014effects}, voltage glitches \cite{11zussa2012investigation, 12zussa2014analysis, 13selmane2011security}, laser \cite{17selmke2016precise, 18dutertre2018laser}, or electromagnetic faults \cite{19dehbaoui2012injection, 20dehbaoui2012electromagnetic}, as well as software, such as Rowhammer \cite{14kim2014flipping}.
It can occur accidentally, for example due to cosmic rays, or deliberately by a mischievous individual; in the latter case, we talk about a fault injection attack.
To deal with fault injection, developers use several countermeasures \cite{dureuil2016,liu2025,Richter-Brockmann2021} or tools \cite{Arlat1990,Arribas2020,1_Segall1995,pensec2024} to protect their programs.
One such tool is called Lazart\cite{Lazart2015}. It is a multi-fault injection tool that mutates the code at the LLVM level in order to simulate the effects of fault injection attacks.

\subsection{Rust}

The Rust language is a recent language defined as a successor to C and C++;
while C/C++ have a lot of memory access bugs such as buffer overflows or data race issues.

Rust has stronger protection over memory access and concurrency over variables; variables are tracked with an element called the borrow checker that tracks the lifespan and ownership, as well as their access.
While the borrow checker protects against bad access of data in the software program. 

The same cannot be said for hardware failure, which may cause some protections to be countered.

\todo[inline]{section 2.2 :Si tu veux plus parler du borrow-checker alors tu peux prendre l'exemple
   de Tommy (échange de pointeurs) du papier FDTC. Tu peux aussi faire réference à la section 4 qui détaille cet exemple ...,ref FDTC ?}

We can take the buffer overflows of a table rust protection as an example.
\texttt{rustc} will add a test on the code at compile time, to identify if you reading in or out the table, and will launch a pannic if it's outside the range of the table.
injected a faute at the moment of the condition can allow the program to read outside the table, and thus counter the protection set by Rust language itself.

An other example, is the change pointer fault. It is explained in the subsection \ref{subsection:Change Pointer fault model}.

\subsection{Contribution}
 
The purpose of this report is to show how Rust differs from C and C++ regarding offensive fault injection.
To answer that, there are two different parts:

\begin{itemize}
\item \textbf{Extend Lazart to support Rust}: discussed in section \ref{Section:Rust ported}, this shows the instructions that have to be added in order to analyze Rust programs without losing attacks. In addition, we also discuss the change pointer fault in this section.
\item \textbf{Porting the FISSC benchmark to Rust}: this point discusses FISSC (Fault Injection and Simulation Secure Collection) \cite{dureuil2016}. We analyze two programs of FISSC in Rust and C and compare them to each other. See section \ref{Section:FISSC}
\end{itemize}

The report is organized as follows. Section \ref{Section:Related_work} presents the Rust programming language, related work concerning Rust and fault injection, and the tool Lazart. Section \ref{Section:Rust ported} describes the extension of Lazart to support the Rust language and the addition of a new attack model for C and Rust languages. Section \ref{Section:CFG} presents a new option of the Lazart tool giving the possibility to generate a control flow graph (CFG). Section \ref{Section:FISSC} shows the comparison made on the FISSC benchmark and a Rust implementation designed to match the C implementation's attack surface. Section \ref{Section:result section} shows results for the FISSC comparison of a Rust implementation, and findings that were produced from control flow graph analysis. And finally, Section \ref{section:Conclusion} summarizes and proposes some perspectives to this work.

%Originally, Lazart only supported C code, but support for Rust is now being implemented. 
%In previous work, a countermeasure benchmark was implemented and evaluated in C \cite{dureuil2016}.

%The purpose of this study is to evaluate Fault injection on Rust code and to compare its results with the ones obtained with C on a representative benchmark.
%Another question addressed in this study, whether there is a way to protect fault injection countermeasures from compiler optimisations, and how much effectiveness remains after the optimisation.

